\chapter{Conjugaison, parité et renversement temporel dans le registre}
\label{chap:cpt-discreto-rutas-rev6}
\index{CPT!symétrie discrète HMT}
\index{action!interaction minimale}
\index{voie!renversement temporel}

La violation CP d'une matrice de mélange et la symétrie CPT ne sont pas la même
affirmation. Avant le théorème CPT d'une théorie quantique des champs, le
registre
HMT possède une symétrie discrète exacte sur les voies. Son contenu est concret :
la conjugaison échange les deux feuillets, la parité réfléchit les directions et
le renversement temporel parcourt le registre en sens inverse. L'action locale
qui compte les rotations et les changements de feuillet est invariante sous les trois opérations.

\section{Voie et action locale}

Une trajectoire de longueur \(L\) est
\[
 \gamma=
 \bigl((i_0,j_0,\ell_0),(d_1,\ldots,d_L)\bigr),
\]
où \(d_t\in\{N,E,S,W\}\) et
\(\ell_t\in\{\Sigma,\Pi\}\). La dynamique TPK détermine les états
\((i_t,j_t,\ell_t)\). Nous définissons
\[
 \operatorname{turn}(\gamma)
 =
 \#\{2\leq t\leq L:d_t\neq d_{t-1}\},
\]
\[
 \operatorname{switch}(\gamma)
 =
 \#\{2\leq t\leq L:\ell_t\neq\ell_{t-1}\}.
\]

\begin{definition}[Action discrète d'interaction minimale]
\label{def:accion-minima-cpt-rev6}
Pour \(\lambda\geq0\),
\[
 \mathcal A_\lambda(\gamma)
 :=
 \operatorname{turn}(\gamma)
 +\lambda\,\operatorname{switch}(\gamma).
\]
De manière équivalente,
\[
 \mathcal A_\lambda(\gamma)
 =
 \sum_{t=2}^{L}
 \left[
 1-\delta_{d_t,d_{t-1}}
 +\lambda(1-\delta_{\ell_t,\ell_{t-1}})
 \right].
\]
\end{definition}

L'action dépend de la variation de la voie, non des noms absolus des
directions ou des feuillets.

\section{Involutions discrètes \(C\), \(P\) et \(T\) sur les voies : conjugaison de feuillet, réflexion spatiale et renversement du registre}

\begin{definition}[Opérateurs discrets \(C,P,T\)]
\label{def:cpt-discreto-rutas-rev6}
Sur une trajectoire complète, on définit :
\begin{enumerate}
\item \(P\), réflexion spatiale :
\[
 (i,j)\longmapsto(i,-j)\pmod9,
 \qquad E\leftrightarrow W,\quad N\mapsto N,\quad S\mapsto S,
\]
sans changer le feuillet ;
\item \(C\), conjugaison de feuillet :
\[
 \Sigma\leftrightarrow\Pi,
\]
sans changer la direction ; dans la carte angulaire, elle correspond à
\(C^*\mapsto-C^*\), qui échange \(q_+\leftrightarrow q_-\) ;
\item \(T\), renversement du registre :
\[
 d_t\longmapsto\operatorname{opp}(d_{L+1-t}),
 \qquad
 \ell_t\longmapsto\ell_{L+1-t},
 \qquad
 (i_t,j_t)\longmapsto(i_{L-t},j_{L-t}).
\]
\end{enumerate}
\end{definition}

Ici \(T\) est une involution de la trajectoire enregistrée. Elle ne s'identifie pas
à l'inversion d'un automate appauvri qui aurait oublié l'état nécessaire
pour reconstruire son prédécesseur.

\begin{theorem}[Invariance \(CPT\) de l'action de voie]
\label{thm:invariancia-cpt-discreta-rev6}
Pour toute trajectoire \(\gamma\) et tout \(\lambda\geq0\),
\[
 \boxed{
 \mathcal A_\lambda(\gamma)
 =
 \mathcal A_\lambda(C\gamma)
 =
 \mathcal A_\lambda(P\gamma)
 =
 \mathcal A_\lambda(T\gamma)
 =
 \mathcal A_\lambda(CPT\,\gamma).}
\]
\end{theorem}

\begin{proof}
\(P\) applique une permutation involutive à l'alphabet des directions. Ainsi
\(d_t=d_{t-1}\) si et seulement si
\(P(d_t)=P(d_{t-1})\), et conserve le nombre de rotations. \(C\) permute
globalement les feuillets, de sorte qu'elle conserve le motif des égalités
\(\ell_t=\ell_{t-1}\) ; elle n'altère pas les directions. \(T\) inverse l'ordre des
deux suites : le nombre de changements entre termes consécutifs est le
même lorsqu'on parcourt un mot dans le sens direct ou inverse. Chaque opération
conserve séparément \(\operatorname{turn}\) et
\(\operatorname{switch}\), donc conserve
\(\mathcal A_\lambda\). La composition la conserve également.
\end{proof}

\section{Parité des lecteurs descendants}

L'échange \(C^*\mapsto-C^*\) permute \(q_+\) et \(q_-\). Par conséquent,
toute expression symétrique dans les deux canaux est \(C\)-paire et toute différence
antisymétrique est \(C\)-impaire. Le
\cref{thm:paridad-bicapa-barbero-rev6} réalise cette règle dans deux objets :
\[
 S_{\rm bi}(-C^*)=S_{\rm bi}(C^*),
 \qquad
 \Gamma_{6\to8}(A,-C^*)
 =-\Gamma_{6\to8}(A,C^*).
\]
La bande de masse du
\cref{thm:principio-conjugacion-masa-banda} réalise la même séparation :
masse paire et charge/orientation impaire.

\section{Relation avec CKM et avec le théorème CPT des champs}

La matrice CKM construite au chapitre suivant possède
\(\delta_{\rm CP}\neq0\) et un invariant de Jarlskog non nul. Cela décrit
une violation CP dans cette représentation de mélange. Cela ne contredit pas le
\cref{thm:invariancia-cpt-discreta-rev6} : ce dernier agit sur
l'ensemble des voies et leur action, sans exiger que chaque lecteur orienté soit pair.

Le théorème CPT de la théorie quantique des champs ajoute un autre domaine et d'autres
hypothèses ---localité relativiste, structure de champs, spectre et
représentation correspondante---. Il est confronté ensuite. Le résultat HMT
démontré ici est la symétrie discrète du registre : il fixe l'objet qu'une
réalisation de champs doit transporter, au lieu de le remplacer par une
référence externe.
